\chapter{Theory} \label{sec:Theory}
It is assumed that the reader has a fundamental understanding of electrical engineering principles. The purpose of this chapter is therefore to present and explain the theoretical concepts that are specifically relevant to this work and may not be familiar to all readers.

\section{Strain}
Strain is a measure of the relative deformation of a material subjected to mechanical loading. It describes how much a body stretches or compresses compared to its original dimensions. For a uniaxial deformation, the engineering strain is defined as:

\begin{align} \label{eq:strain definition}
    \varepsilon=\frac{\Delta L}{L_{0}}
\end{align}

where $\Delta L$ is the change in length and $L_0$ is the original length of the material. For small deformations, this definition provides an accurate and widely used description of material behavior.

During assembly of the bolted joint, strain occurs both in the bolt stretching and the nut compressing, see \cref{fig:strain}. This deformation is directly related to the preload in the bolt, and a reduction in preload will therefore result in a corresponding change in strain. This relationship enables strain to be used as an indicator for monitoring the mechanical state of the joint.

\begin{figure}[H]
\centering
\includestandalone[mode=buildmissing, width=\textwidth]{Standalone/Drawings/strain-tikz}
    \caption{Illustration of strain in a bolted joint. The bolt undergoes tensile stretching ($\Delta L_b$), while the two nuts experience a corresponding compressive deformation ($\Delta L_n$), due to the applied preload.}
    \label{fig:strain}
\end{figure}

% The strain requirement for the system considered in this work is $800\hspace{0.07cm} \text{\textmu}$. For a characteristic length of $\SI{30}{mm}$ along the nut surface, this corresponds to a minimum detectable stretch of:

% \begin{equation} \label{eq:strain requirement}
% \begin{aligned} 
%     \varepsilon = 800\hspace{0.07cm} \text{\textmu} &=\frac{\SI{800}{\micro \meter}}{\SI{1}{\meter}} \\
%     &= \frac{\SI{24}{\micro \meter}}{\SI{30}{\milli \meter}}
% \end{aligned}
% \end{equation}

\pagebreak
\section{Constriction Resistance}
Constriction resistance refers to the additional electrical resistance that arises at the interface between two conductive bodies due to the discrete nature of their mechanical contact. Although surfaces may appear smooth at the macroscopic scale, they are inherently rough when observed microscopically, consisting of asperities and voids, see \cref{fig:CR}. As a result, electrical contact is established only at a limited number of contact points rather than across the entire apparent contact area.
At these asperity contacts, the electrical current is forced to converge and pass through highly localized regions, leading to a constriction of the current flow. This geometric restriction increases the effective resistance compared to an ideal continuous conductor. The total constriction resistance is therefore governed by both the number of asperities and their individual contact sizes, which is strongly influenced by mechanical loading. 
As the load between the conductive surfaces increases, additional asperities come into contact and existing contact areas increase. This results in a larger effective region and a corresponding decrease in constriction resistance. Conversely, at low contact forces, only a few asperities contribute to conduction, leading to higher resistance.
This dependency enables constriction resistance to be exploited as a sensing mechanism, where mechanical stimuli are transduced into measurable changes in electrical resistance \cite{G.O.A.T}.

\begin{figure}[H]
\centering
\includestandalone[mode=buildmissing]{Standalone/Drawings/CR-tikz}
    \caption{Illustration of microscopic view between two surfaces creating contact with multiple asperities. Top figure with low force applied, and bottom with high force applied.}
    \label{fig:CR}
\end{figure}